\chapter{المجسّمات والحجوم}\label{ch:g9:solids}

تبدأ هندسة الفضاء بأسرة صغيرة من \omterm{def:g5:solids:def}{المجسّمات} --- \omterm{def:g7:areas:prism}{المواشير}،
والأسطوانات، والأهرام، \omterm{def:g8:solids:def}{والمخاريط}، والكرات --- وبسؤالين: كم
تسع (\omterm{def:g6:measure:volume}{الحجم})، وماذا ترى حين تقطعها (المقاطع)؟ وينتهي الفصل
بواقعة بالغة الأهمية العملية: تحجيم \omterm{def:g5:solids:def}{مجسّم} بالمعامل $k$ يضرب
حجمه في $k^3$.

\section{المجسّمات الكلاسيكية وحجومها}

\begin{theorem}[صيغ الحجوم]\label{thm:g9:solids:volumes}
لنكتب $B$ \omterm{def:g6:measure:area}{لمساحة} القاعدة و $h$ للاِرتفاع (وهو المسافة بين
القاعدة والوجه المقابل أو الرأس).
\begin{center}
\renewcommand{\arraystretch}{1.5}
\begin{tabular}{l|l}
\omterm{def:g5:solids:def}{المجسّم} & \omterm{def:g6:measure:volume}{الحجم} \\
\hline
\omterm{def:g7:areas:prism}{الموشور}، \omterm{def:g7:areas:prism}{والأسطوانة} & $V = B \times h$ \\[5pt]
\omterm{def:g8:solids:def}{الهرم}، \omterm{def:g8:solids:def}{والمخروط} & $V = \dfrac13\, B \times h$ \\[5pt]
كرة نصف قطرها $r$ & $V = \dfrac43\,\pi r^3$ \\
\end{tabular}
\end{center}
ومن أجل \omterm{def:g7:areas:prism}{أسطوانة} \omterm{def:g8:solids:def}{ومخروط} نصف قطرهما $r$، تكون \omterm{def:g6:measure:area}{مساحة} القاعدة
$B = \pi r^2$؛ \omterm{def:g6:measure:area}{ومساحة} سطح الكرة $4\pi r^2$.
\end{theorem}
\admitted

\begin{omfigure}
\begin{tikzpicture}[scale=0.62]
  % cylinder
  \draw[thick] (0,0) ellipse (1 and 0.3);
  \draw[thick] (-1,0) -- (-1,2.6);
  \draw[thick] (1,0) -- (1,2.6);
  \draw[thick] (0,2.6) ellipse (1 and 0.3);
  \draw[dashed] (0,0) -- (0,2.6);
  \node[below, font=\small] at (0,-0.4) {أسطوانة};
  % cone
  \begin{scope}[xshift=4.2cm]
  \draw[thick] (0,0) ellipse (1 and 0.3);
  \draw[thick] (-1,0) -- (0,2.8);
  \draw[thick] (1,0) -- (0,2.8);
  \draw[dashed] (0,0) -- (0,2.8);
  \node[below, font=\small] at (0,-0.4) {مخروط};
  \end{scope}
  % pyramid
  \begin{scope}[xshift=8.4cm]
  \draw[thick] (-1.1,0) -- (0.9,0) -- (1.5,0.55) -- (-0.5,0.55) -- cycle;
  \draw[thick] (-1.1,0) -- (0.2,2.8);
  \draw[thick] (0.9,0) -- (0.2,2.8);
  \draw[thick] (1.5,0.55) -- (0.2,2.8);
  \draw[dashed] (-0.5,0.55) -- (0.2,2.8);
  \node[below, font=\small] at (0.2,-0.4) {هرم};
  \end{scope}
  % sphere
  \begin{scope}[xshift=12.8cm]
  \draw[thick] (0,1.3) circle (1.3);
  \draw[dashed] (0,1.3) ellipse (1.3 and 0.4);
  \node[below, font=\small] at (0,-0.4) {كرة};
  \end{scope}
\end{tikzpicture}

{\small \omterm{def:g5:solids:def}{مجسّمات} هذا الفصل. ويضمّ \omterm{def:g6:measure:volume}{حجم} كل واحد منها \omterm{def:g6:measure:area}{مساحة}
قاعدة واِرتفاعًا --- إلا الكرة، فلا تحتاج سوى نصف قطرها.}
\end{omfigure}

\begin{example}\label{ex:g9:solids:volumes}
علبة أسطوانية نصف قطرها $4$ سم واِرتفاعها $10$ سم:
\[
V = \pi r^2 h = \pi \times 16 \times 10 = 160\pi \approx 503 \text{ سم}^3
\approx 0.5 \text{ لتر}.
\]
\omterm{def:g8:solids:def}{ومخروط} له القاعدة والاِرتفاع نفسهما يسع ثلث ذلك:
$\frac{160\pi}{3} \approx 168$ سم$^3$. وكرة نصف قطرها $4$ سم:
$V = \frac43 \pi \times 64 = \frac{256\pi}{3} \approx 268$ سم$^3$.
\end{example}

\begin{example}[هرم خطوةً خطوة]\label{ex:g9:solids:pyramid}
\omterm{def:g8:solids:def}{هرم} قاعدته مربّعة ضلعها $6$ م واِرتفاعه $5$ م.
\begin{enumerate}
  \item \omterm{def:g6:measure:area}{مساحة} القاعدة: $B = 6^2 = 36$ م$^2$.
  \item \omterm{def:g6:measure:volume}{والحجم}: $V = \frac13 \times 36 \times 5 = 60$ م$^3$.
\end{enumerate}
واِنتبه للوحدات: لو أُعطي الضلع بالسنتيمتر والاِرتفاع بالمتر،
لوجب تحويل أحدهما أولًا.
\end{example}

\section{المقاطع بمستويات}

\begin{proposition}[مقاطع المجسّمات الكلاسيكية]\label{prop:g9:solids:sections}
قطع \omterm{def:g5:solids:def}{مجسّم} بمستوٍ يعطي شكلًا مستويًا، هو
\emph{المقطع}\index{مقطع}:
\begin{itemize}
  \item \omterm{def:g7:areas:prism}{الموشور} أو \omterm{def:g7:areas:prism}{الأسطوانة} المقطوع \emph{\omterm{def:g6:lines:perp}{موازيًا} لقاعدته}
        يعطي نسخة من القاعدة، عند كل اِرتفاع؛
  \item \omterm{def:g8:solids:def}{والهرم} أو \omterm{def:g8:solids:def}{المخروط} المقطوع \omterm{def:g6:lines:perp}{موازيًا} لقاعدته يعطي
        \emph{تصغيرًا} للقاعدة: فعلى بعد $d$ من الرأس، يكون
        \omterm{def:g9:thales:scaling}{معامل التحجيم} $k = \frac{d}{h}$؛
  \item والكرة ذات \omterm{def:g3:shapes:shapes}{نصف القطر} $r$ المقطوعة بمستوٍ على بعد
        $d < r$ من المركز تعطي دائرة نصف قطرها
        $\sqrt{r^2 - d^2}$ (حسب مبرهنة فيثاغورس).
\end{itemize}
\end{proposition}
\admitted

\begin{omfigure}
\begin{tikzpicture}[scale=0.75]
  % cone with section
  \draw[thick] (0,0) ellipse (1.5 and 0.45);
  \draw[thick] (-1.5,0) -- (0,4.0);
  \draw[thick] (1.5,0) -- (0,4.0);
  \draw[omThm, thick, fill=omThm!12] (0,2.4) ellipse (0.6 and 0.18);
  \draw[dashed] (0,0) -- (0,4.0);
  \node[right, font=\small] at (0.1,4.0) {الرأس};
  \draw[Stealth-Stealth, thin] (2.1,0) -- (2.1,4.0);
  \node[right, font=\small] at (2.1,2.0) {$h$};
  \draw[Stealth-Stealth, thin, omThm] (0.9,2.4) -- (0.9,4.0);
  \node[right, font=\small, omThm] at (0.92,3.2) {$d$};
  % sphere with section
  \begin{scope}[xshift=7.6cm]
  \draw[thick] (0,2) circle (1.8);
  \draw[omThm, thick, fill=omThm!12] (0,2.9) ellipse (1.56 and 0.4);
  \fill (0,2) circle (1.5pt) node[below, font=\small] {$O$};
  \draw[thin] (0,2) -- (0,2.9) node[midway, left, font=\small] {$d$};
  \draw[thin] (0,2.9) -- (1.56,2.9)
    node[midway, above, font=\small] {$\sqrt{r^2 - d^2}$};
  \draw[thin] (0,2) -- (-1.27,0.72)
    node[midway, below right=-3pt, font=\small] {$r$};
  \end{scope}
\end{tikzpicture}

{\small على اليسار: قطع \omterm{def:g8:solids:def}{مخروط} \omterm{def:g6:lines:perp}{موازيًا} لقاعدته على بعد $d$ من
الرأس يعطي قرصًا محجَّمًا بالمعامل $k = \frac dh$. وعلى اليمين:
مستوٍ على بعد $d$ من مركز كرة يقطعها في دائرة نصف قطرها
$\sqrt{r^2 - d^2}$.}
\end{omfigure}

\begin{example}\label{ex:g9:solids:section}
\omterm{def:g8:solids:def}{مخروط} \omterm{def:g3:shapes:shapes}{نصف قطر} قاعدته $6$ واِرتفاعه $9$. \omterm{prop:g9:solids:sections}{والمقطع} بمستوٍ
\omterm{def:g6:lines:perp}{موازٍ} للقاعدة، على بعد $3$ من الرأس، قرص محجَّم بالمعامل
$k = \frac39 = \frac13$: فنصف قطره $6 \times \frac13 = 2$.

وكرة نصف قطرها $5$ تُقطع بمستوٍ على بعد $3$ من مركزها:
\omterm{prop:g9:solids:sections}{فالمقطع} دائرة نصف قطرها $\sqrt{25 - 9} = 4$.
\end{example}

\section{تحجيم المجسّمات}

\begin{theorem}[أثر التحجيم في الأطوال والمساحات والحجوم]\label{thm:g9:solids:scaling}
حين يُكبَّر \omterm{def:g5:solids:def}{مجسّم} أو يُصغَّر \omterm{def:g9:thales:scaling}{بمعامل التحجيم} $k > 0$:
\begin{itemize}
  \item تُضرب كل الأطوال في $k$؛
  \item وتُضرب كل المساحات في $k^2$؛
  \item وتُضرب كل الحجوم في $k^3$.
\end{itemize}
\end{theorem}
\admitted

\begin{example}\label{ex:g9:solids:scaling}
نموذج سيارة \omterm{def:g7:prop:scale}{بسلّم} $\frac{1}{18}$: فأطواله هي أطوال السيارة
الحقيقية مضروبة في $k = \frac{1}{18}$، وسطحه المطليّ مضروب في
$k^2 = \frac{1}{324}$، وحجمه في $k^3 = \frac{1}{5832}$.

ومضاعفة \omterm{def:g3:shapes:shapes}{نصف قطر} كرة ($k = 2$) تضرب حجمها في
$2^3 = 8$: تحقّق على الصيغة،
$\frac43\pi (2r)^3 = \frac43 \pi \times 8r^3$.
\end{example}

\begin{example}[المخروط الناقص]\label{ex:g9:solids:truncated}
\omterm{def:g8:solids:def}{مخروط} اِرتفاعه $9$ \omterm{def:g3:shapes:shapes}{ونصف قطر} قاعدته $6$ (وحجمه
$\frac13 \pi \times 36 \times 9 = 108\pi$) يُقطع على بعد $3$ من
الرأس، ويُزال \omterm{def:g8:solids:def}{المخروط} الصغير فوق \omterm{def:g6:lines:objects}{القطع}. \omterm{def:g8:solids:def}{والمخروط} الصغير هو
التصغير بالمعامل $k = \frac13$، فحجمه
$108\pi \times \left(\frac13\right)^3 = 4\pi$. \omterm{def:g5:solids:def}{والمجسّم} \omterm{def:g3:division:remainder}{الباقي}
(\emph{\omterm{def:g8:solids:def}{مخروط} ناقص}) حجمه $108\pi - 4\pi = 104\pi \approx 327$.
\end{example}

\section{تمارين}

\begin{exercise}[$\star$]\label{exo:g9:solids:1}
اِحسب الحجوم: علبة (\omterm{def:g7:areas:prism}{موشور} قائم) أبعادها $4 \times 5
\times 12$؛ \omterm{def:g7:areas:prism}{وأسطوانة} نصف قطرها $3$ واِرتفاعها $7$ (القيمة
المضبوطة بالرمز $\pi$، ثم مقرّبة إلى الوحدة).
\end{exercise}

\begin{exercise}[$\star$]\label{exo:g9:solids:2}
اِحسب \omterm{def:g6:measure:volume}{حجم} \omterm{def:g8:solids:def}{مخروط} نصف قطره $5$ واِرتفاعه $12$، \omterm{def:g8:solids:def}{وهرم}
قاعدته مستطيلة $8 \times 3$ واِرتفاعه $10$.
\end{exercise}

\begin{exercise}[$\star$]\label{exo:g9:solids:3}
اِحسب \omterm{def:g6:measure:volume}{حجم} \omterm{def:g6:measure:area}{ومساحة} سطح كرة نصف قطرها $6$ (القيمتين
المضبوطتين بالرمز $\pi$).
\end{exercise}

\begin{exercise}[$\star$]\label{exo:g9:solids:4}
كرة نصف قطرها $10$ تُقطع بمستوٍ على بعد $8$ من مركزها. فما
\omterm{def:g3:shapes:shapes}{نصف قطر} دائرة \omterm{prop:g9:solids:sections}{المقطع}؟
\end{exercise}

\begin{exercise}[$\star\star$]\label{exo:g9:solids:5}
\omterm{def:g8:solids:def}{مخروط} \omterm{def:g3:shapes:shapes}{نصف قطر} قاعدته $8$ واِرتفاعه $12$. ويقطعه مستوٍ \omterm{def:g6:lines:perp}{موازٍ}
للقاعدة على بعد $9$ من الرأس. اِحسب \omterm{def:g3:shapes:shapes}{نصف قطر} \omterm{prop:g9:solids:sections}{المقطع}، ثم \omterm{def:g6:measure:volume}{حجم}
\omterm{def:g8:solids:def}{المخروط} الصغير فوق \omterm{def:g6:lines:objects}{القطع}.
\end{exercise}

\begin{exercise}[$\star\star$]\label{exo:g9:solids:6}
\omterm{def:g8:powers:def}{كأس} أسطوانية نصف قطرها الداخلي $3$ سم فيها ماء إلى اِرتفاع
$10$ سم. وتُغمر فيها كرة نصف قطرها $2$ سم غمرًا تامًّا. فكم
يرتفع مستوى الماء؟ (\omterm{def:g6:measure:volume}{فالحجم} المضاف ينتشر على \omterm{def:g6:measure:area}{مساحة} قاعدة
الكأس؛ أعطِ الاِرتفاع المضبوط، ثم قرّب إلى المليمتر.)
\end{exercise}

\begin{exercise}[$\star\star$]\label{exo:g9:solids:7}
وصفة تملأ قالبًا كرويًّا نصف قطره $6$ سم. وليس لديك سوى قالب
كرويّ نصف قطره $3$ سم. فكم قالبًا صغيرًا تستطيع أن تملأ؟
(أجِب من غير حساب أيّ من الحجمين صراحةً.)
\end{exercise}

\begin{exercise}[$\star\star$]\label{exo:g9:solids:8}
\omterm{def:g8:solids:def}{الهرم} الأكبر بالجيزة قاعدته مربّعة ضلعها نحو $230$ م
واِرتفاعه نحو $147$ م. قدّر حجمه، وعبّر عنه بملايين الأمتار
المكعّبة.
\end{exercise}

\begin{exercise}[$\star\star\star$]\label{exo:g9:solids:9}
قمع مخروطي نصف قطره $6$ سم واِرتفاعه $12$ سم، ورأسه إلى
الأسفل، يُملأ ماءً إلى نصف \emph{اِرتفاعه} (مقيسًا من
الرأس).
\begin{enumerate}
  \item فأيّ \omterm{def:g9:fractions:fraction}{كسر} من \omterm{def:g6:measure:volume}{حجم} القمع مملوء؟
  \item وإذا مُلئ بدلًا من ذلك بنصف \emph{حجمه} ماءً، بيّن أنّ
        اِرتفاع الماء $d$ يحقّق
        $\left(\frac{d}{12}\right)^3 = \frac12$، وأعطِ $d$ إلى
        المليمتر ($\sqrt[3]{0.5} \approx 0.7937$).
\end{enumerate}
\end{exercise}

\section{مسألة: شاهد قبر أرخميدس}

\begin{problem}[{مسألة نهاية الأسبوع --- الكرة ثلثا أسطوانتها
(مرّتين)، ورصّة 1:2:3، وقانون المربّع والمكعّب الذي يمنع
العمالقة}]
\label{pb:g9:solids:1}
برهن أرخميدس على مبرهنات كثيرة، لكنّ واحدة منها أفخرته حتى طلب
أن يُنقش \emph{شكلها} على قبره: كرة محشورة في أضيق \omterm{def:g7:areas:prism}{أسطوانة}
لها. وبعد قرن، تعرّف الكاتب الروماني شيشرون، وهو يفتّش في
العلّيق قرب سيراكوزا، على القبر «بالكرة \omterm{def:g7:areas:prism}{والأسطوانة}». وتستعيد
هذه المسألة ما يرمز إليه النقش --- معجزة ثلثين مزدوجة --- ثم
تتبع الحجوم والسطوح (\cref{thm:g9:solids:volumes}
و\cref{thm:g9:solids:scaling}) إلى قانون يحكم العمالقة والنمل
والكواكب التي تبرد.

\medskip
\noindent\textbf{الجزء الأول --- فكّ رموز النقش.}
كرة نصف قطرها $r$ تجلس بالضبط داخل \omterm{def:g7:areas:prism}{أسطوانة}: \omterm{def:g3:shapes:shapes}{بنصف القطر}
نفسه، واِرتفاعها $2r$.
\begin{enumerate}
  \item اِحسب \omterm{def:g6:measure:volume}{حجم} \omterm{def:g7:areas:prism}{الأسطوانة} بدلالة $r$.
  \item واِحسب نسبة \omterm{def:g6:measure:volume}{حجم} الكرة إلى \omterm{def:g6:measure:volume}{حجم} \omterm{def:g7:areas:prism}{الأسطوانة}. ولماذا قد
        يكون أرخميدس أحبّ أنّ الجواب لا يحوي $\pi$ ولا $r$؟
  \item والآن السطوح: قارن \omterm{def:g6:measure:area}{مساحة} سطح الكرة بسطح \omterm{def:g7:areas:prism}{الأسطوانة}
        \emph{الكلّي} (الجانب مضافًا إليه الغطاءان). فأيّ نسبة
        تظهر --- من جديد؟
  \item والفضاء الفارغ بين الكرة \omterm{def:g7:areas:prism}{والأسطوانة} حجمه
        $2\pi r^3 - \frac43\pi r^3$. بيّن أنّ هذا \omterm{def:g3:division:remainder}{الباقي} يساوي
        بالضبط \omterm{def:g6:measure:volume}{حجم} \emph{\omterm{def:g8:solids:def}{مخروطين}} نصف قطرهما
        $r$ واِرتفاعهما $r$.
  \item وكرة سلّة نصف قطرها $12$ سم تُباع في أضيق علبة
        أسطوانية. اِحسب الحجمين (إلى أقرب
        $100$ سم$^3$) والنسبة المئوية للجزء الفارغ من
        العلبة.
\end{enumerate}

\medskip
\noindent\textbf{الجزء الثاني --- الأوعية والقوالب والأقمار.}
\begin{enumerate}[resume]
  \item وعاء نصف كروي نصف قطره $10$ سم. اِحسب سعته،
        بوحدة سم$^3$ وباللترات (إلى الديسيلتر).
  \item ورصّة 1:2:3: \omterm{def:g8:solids:def}{مخروط}، ونصف كرة، \omterm{def:g7:areas:prism}{وأسطوانة}،
        وكلّها نصف قطرها $10$ سم واِرتفاعها $10$ سم. اِحسب
        الحجوم الثلاثة وتحقّق من النسبة الأسطورية
        $1 : 2 : 3$. (وكان أرخميدس ليقدّر هذه أيضًا.)
  \item وكرة شوكولاتة نصف قطرها $3$ سم تُذاب في قالب
        أسطواني نصف قطره $3$ سم. فأيّ اِرتفاع تبلغه
        الشوكولاتة؟ (كسرًا مضبوطًا، ثم إلى
        المليمتر.)
  \item \omterm{def:g3:shapes:shapes}{ونصف قطر} الأرض نحو $6\,371$ كلم، \omterm{def:g3:shapes:shapes}{ونصف قطر} القمر نحو
        $1\,737$ كلم. اِحسب نسبة نصفَي \omterm{def:g4:geometry:circle}{القطر}، ثم --- حسب
        \cref{thm:g9:solids:scaling} --- نسبة الحجوم: فكم
        قمرًا يسع في أرض جوفاء، حجمًا؟
  \item ومن المبرهنة نفسها: ما نسبة \emph{مساحتَي السطح} بين
        الأرض والقمر؟ وفي العموم، حين يتضاعف \omterm{def:g3:shapes:shapes}{نصف قطر} بالون،
        ماذا يحدث لحجمه ولسطحه؟ اُذكر \emph{قانون المربّع
        والمكعّب}: فحين يتحجّم شكل إلى أعلى، تسبق الحجوم
        السطوح.
\end{enumerate}

\medskip
\noindent\textbf{الجزء الثالث --- قانون المربّع والمكعّب يحكم
العالم.}
\begin{enumerate}[resume]
  \item حجّة غاليليو ضدّ العمالقة: تخيّل إنسانًا محجَّمًا
        بالمعامل $10$، بالنسب نفسها. فبأيّ معامل ينمو الوزن
        (والوزن يتبع \omterm{def:g6:measure:volume}{الحجم})؟ وبأيّ معامل ينمو مقطع العظام
        (وهو \omterm{def:g6:measure:area}{مساحة})؟ وبأيّ معامل إذن ينمو \emph{الضغط} على كل
        سنتيمتر مربّع من العظم --- وماذا يحدث للعملاق؟
  \item والقانون نفسه معكوسًا يفسّر بطولات النمل: فالقوّة تتبع
        مقطع العضلة (وهو \omterm{def:g6:measure:area}{مساحة})، والوزن يتبع \omterm{def:g6:measure:volume}{الحجم}. ومن أجل
        حيوان أصغر $100$ مرّة في كل اِتّجاه، بأيّ معاملين ينكمش
        الوزن والقوّة، وبأيّ معامل تتحسّن نسبة القوّة إلى الوزن؟
  \item والحرارة ينتجها \omterm{def:g6:measure:volume}{حجم} الجسم وتُفقد عبر سطحه. بيّن أنّ
        نسبة السطح إلى \omterm{def:g6:measure:volume}{الحجم} في الكرة هي $\frac{3}{r}$،
        واِحسبها من أجل $r = 1$ و $r = 2$. فأيّهما يبرد أسرع،
        الفأر أم الدبّ --- ولماذا يحتاج الرضّع قبّعات في
        الشتاء؟
  \item وبرتقالتان كرويتان نصفا قطريهما $4$ سم و $5$ سم
        تُباعان بالثمن نفسه. اِحسب نسبة حجميهما: فكم من
        البرتقال تعطي الكبيرة أكثر مقابل المال نفسه؟
  \item والختام: صِف بدقّة ما يرمز إليه شكل شاهد القبر ---
        «الثلثان» في السؤالين 2 و 3 --- وأجِب شيشرون في جملة
        واحدة: لماذا يختار عالم رياضيات، على كل فتوحات عبقريته
        الهندسية، كرةً في \omterm{def:g7:areas:prism}{أسطوانة}؟
\end{enumerate}
\end{problem}
